% lemma1_borel_gcf.tex — Lemma 1: k-Shift Borel Regularization Identity
% Self-contained snippet; paste into your paper body or % lemma1_borel_gcf.tex — Lemma 1: k-Shift Borel Regularization Identity
% Self-contained snippet; paste into your paper body or % lemma1_borel_gcf.tex — Lemma 1: k-Shift Borel Regularization Identity
% Self-contained snippet; paste into your paper body or % lemma1_borel_gcf.tex — Lemma 1: k-Shift Borel Regularization Identity
% Self-contained snippet; paste into your paper body or \input{lemma1_borel_gcf}
%
% Required packages (add to preamble if not already present):
%   \usepackage{amsmath,amssymb,amsthm}
%   \newtheorem{lemma}{Lemma}

\begin{lemma}[$k$-Shift Borel Regularization Identity]\label{lem:borel-gcf}
Let the generalized continued fraction (GCF) be defined by partial numerators
$a_n = -n!$ and constant partial denominators $b_n = k$ ($k > 0$),
with~$b_0 = k$.  The formal continued fraction is classically divergent
(by the Stern--Stolz theorem, since $|a_n/b_n| \to \infty$),
but its Borel regularization exists and equals
\[
  V(k)
  \;=\;\int_{0}^{\infty}\frac{k\,e^{-t}}{k+t}\,dt
  \;=\; k\,e^{k}\,E_1(k),
\]
where\/ $E_1(z) = \displaystyle\int_z^{\infty}\frac{e^{-t}}{t}\,dt$
is the exponential integral.
For $k=2$ we obtain $V(2) = 2e^{2}E_1(2) \approx 0.72266\ldots$
\end{lemma}

\begin{proof}
The GCF with $a_n = -n!$ and $b_n = 1$ is formally associated
(via Euler's connection between continued fractions and power series)
to the alternating factorial series $f(x) = \sum_{n \ge 0} (-1)^n\,n!\,x^n$.
Its Borel transform cancels the factorial growth:
\[
  \mathcal{B}[f](t)
  \;=\;\sum_{n \ge 0}\frac{(-1)^n\,n!}{n!}\,t^n
  \;=\;\sum_{n \ge 0}(-t)^n
  \;=\;\frac{1}{1+t},
\]
in the canonical $k = 1$ normalization.

Multiplying all partial denominators by~$k$ is an equivalence
transformation of the continued fraction (Perron~\cite{Perron1957}).
In the Borel plane this rescales the Stieltjes kernel:
the substitution $t \mapsto t/k$ moves the simple pole from
$t = -1$ to $t = -k$, giving
\[
  \mathcal{B}_k(t) \;=\; \frac{k}{k+t}.
\]

Laplace (Borel) resummation along $[0,\infty)$ then yields
\[
  V(k)
  \;=\;\int_{0}^{\infty} e^{-t}\,\mathcal{B}_k(t)\,dt
  \;=\;\int_{0}^{\infty}\frac{k\,e^{-t}}{k+t}\,dt.
\]
The standard identity
$\displaystyle\int_0^\infty \frac{e^{-t}}{x+t}\,dt = e^{x}\,E_1(x)$
(valid for $\operatorname{Re} x > 0$; see~\cite[Ch.\,5]{DLMF})
with $x = k$ completes the proof:
\[
  V(k) = k\,e^{k}\,E_1(k). \qedhere
\]
\end{proof}

\paragraph{Numeric verification.}
A backward-recurrence computation of the GCF convergents at depth
$N = 200$ with $120$-digit precision matches the closed form to all
displayed digits.  For $k = 1,2,3$:
\[
\begin{array}{c|l|l}
k & V(k) = k\,e^k\,E_1(k) & \text{50-digit verification} \\
\hline
1 & 0.59634736232319407434\ldots & \checkmark \\
2 & 0.72265948505688951953\ldots & \checkmark \\
3 & 0.78606692861534416385\ldots & \checkmark
\end{array}
\]


% ── Bibliography entries (add to your .bib or \begin{thebibliography}) ──
%
% @book{Perron1957,
%   author    = {Perron, Oskar},
%   title     = {Die Lehre von den Kettenbr\"uchen},
%   edition   = {3rd},
%   publisher = {Teubner},
%   address   = {Stuttgart},
%   year      = {1957},
%   note      = {Equivalence transformations: \S\S\,1--4}
% }
%
% @book{Hardy1949,
%   author    = {Hardy, G. H.},
%   title     = {Divergent Series},
%   publisher = {Clarendon Press},
%   address   = {Oxford},
%   year      = {1949},
%   note      = {Borel summability: Ch.\,VIII}
% }
%
% @misc{DLMF,
%   title     = {{NIST} Digital Library of Mathematical Functions},
%   howpublished = {\url{https://dlmf.nist.gov/}},
%   note      = {Exponential integral $E_1$: \S\,6.2}
% }
%
% @book{Stieltjes1894,
%   author    = {Stieltjes, Thomas Jan},
%   title     = {Recherches sur les fractions continues},
%   journal   = {Annales de la Facult\'e des Sciences de Toulouse},
%   volume    = {8},
%   pages     = {J1--J122},
%   year      = {1894},
%   note      = {Stieltjes transform and moment representation}
% }
%
% @book{LorentzenWaadeland2008,
%   author    = {Lorentzen, Lisa and Waadeland, Haakon},
%   title     = {Continued Fractions: Convergence Theory},
%   edition   = {2nd},
%   publisher = {Atlantis Press},
%   year      = {2008},
%   note      = {Stern--Stolz theorem: Theorem~3.21}
% }

%
% Required packages (add to preamble if not already present):
%   \usepackage{amsmath,amssymb,amsthm}
%   \newtheorem{lemma}{Lemma}

\begin{lemma}[$k$-Shift Borel Regularization Identity]\label{lem:borel-gcf}
Let the generalized continued fraction (GCF) be defined by partial numerators
$a_n = -n!$ and constant partial denominators $b_n = k$ ($k > 0$),
with~$b_0 = k$.  The formal continued fraction is classically divergent
(by the Stern--Stolz theorem, since $|a_n/b_n| \to \infty$),
but its Borel regularization exists and equals
\[
  V(k)
  \;=\;\int_{0}^{\infty}\frac{k\,e^{-t}}{k+t}\,dt
  \;=\; k\,e^{k}\,E_1(k),
\]
where\/ $E_1(z) = \displaystyle\int_z^{\infty}\frac{e^{-t}}{t}\,dt$
is the exponential integral.
For $k=2$ we obtain $V(2) = 2e^{2}E_1(2) \approx 0.72266\ldots$
\end{lemma}

\begin{proof}
The GCF with $a_n = -n!$ and $b_n = 1$ is formally associated
(via Euler's connection between continued fractions and power series)
to the alternating factorial series $f(x) = \sum_{n \ge 0} (-1)^n\,n!\,x^n$.
Its Borel transform cancels the factorial growth:
\[
  \mathcal{B}[f](t)
  \;=\;\sum_{n \ge 0}\frac{(-1)^n\,n!}{n!}\,t^n
  \;=\;\sum_{n \ge 0}(-t)^n
  \;=\;\frac{1}{1+t},
\]
in the canonical $k = 1$ normalization.

Multiplying all partial denominators by~$k$ is an equivalence
transformation of the continued fraction (Perron~\cite{Perron1957}).
In the Borel plane this rescales the Stieltjes kernel:
the substitution $t \mapsto t/k$ moves the simple pole from
$t = -1$ to $t = -k$, giving
\[
  \mathcal{B}_k(t) \;=\; \frac{k}{k+t}.
\]

Laplace (Borel) resummation along $[0,\infty)$ then yields
\[
  V(k)
  \;=\;\int_{0}^{\infty} e^{-t}\,\mathcal{B}_k(t)\,dt
  \;=\;\int_{0}^{\infty}\frac{k\,e^{-t}}{k+t}\,dt.
\]
The standard identity
$\displaystyle\int_0^\infty \frac{e^{-t}}{x+t}\,dt = e^{x}\,E_1(x)$
(valid for $\operatorname{Re} x > 0$; see~\cite[Ch.\,5]{DLMF})
with $x = k$ completes the proof:
\[
  V(k) = k\,e^{k}\,E_1(k). \qedhere
\]
\end{proof}

\paragraph{Numeric verification.}
A backward-recurrence computation of the GCF convergents at depth
$N = 200$ with $120$-digit precision matches the closed form to all
displayed digits.  For $k = 1,2,3$:
\[
\begin{array}{c|l|l}
k & V(k) = k\,e^k\,E_1(k) & \text{50-digit verification} \\
\hline
1 & 0.59634736232319407434\ldots & \checkmark \\
2 & 0.72265948505688951953\ldots & \checkmark \\
3 & 0.78606692861534416385\ldots & \checkmark
\end{array}
\]


% ── Bibliography entries (add to your .bib or \begin{thebibliography}) ──
%
% @book{Perron1957,
%   author    = {Perron, Oskar},
%   title     = {Die Lehre von den Kettenbr\"uchen},
%   edition   = {3rd},
%   publisher = {Teubner},
%   address   = {Stuttgart},
%   year      = {1957},
%   note      = {Equivalence transformations: \S\S\,1--4}
% }
%
% @book{Hardy1949,
%   author    = {Hardy, G. H.},
%   title     = {Divergent Series},
%   publisher = {Clarendon Press},
%   address   = {Oxford},
%   year      = {1949},
%   note      = {Borel summability: Ch.\,VIII}
% }
%
% @misc{DLMF,
%   title     = {{NIST} Digital Library of Mathematical Functions},
%   howpublished = {\url{https://dlmf.nist.gov/}},
%   note      = {Exponential integral $E_1$: \S\,6.2}
% }
%
% @book{Stieltjes1894,
%   author    = {Stieltjes, Thomas Jan},
%   title     = {Recherches sur les fractions continues},
%   journal   = {Annales de la Facult\'e des Sciences de Toulouse},
%   volume    = {8},
%   pages     = {J1--J122},
%   year      = {1894},
%   note      = {Stieltjes transform and moment representation}
% }
%
% @book{LorentzenWaadeland2008,
%   author    = {Lorentzen, Lisa and Waadeland, Haakon},
%   title     = {Continued Fractions: Convergence Theory},
%   edition   = {2nd},
%   publisher = {Atlantis Press},
%   year      = {2008},
%   note      = {Stern--Stolz theorem: Theorem~3.21}
% }

%
% Required packages (add to preamble if not already present):
%   \usepackage{amsmath,amssymb,amsthm}
%   \newtheorem{lemma}{Lemma}

\begin{lemma}[$k$-Shift Borel Regularization Identity]\label{lem:borel-gcf}
Let the generalized continued fraction (GCF) be defined by partial numerators
$a_n = -n!$ and constant partial denominators $b_n = k$ ($k > 0$),
with~$b_0 = k$.  The formal continued fraction is classically divergent
(by the Stern--Stolz theorem, since $|a_n/b_n| \to \infty$),
but its Borel regularization exists and equals
\[
  V(k)
  \;=\;\int_{0}^{\infty}\frac{k\,e^{-t}}{k+t}\,dt
  \;=\; k\,e^{k}\,E_1(k),
\]
where\/ $E_1(z) = \displaystyle\int_z^{\infty}\frac{e^{-t}}{t}\,dt$
is the exponential integral.
For $k=2$ we obtain $V(2) = 2e^{2}E_1(2) \approx 0.72266\ldots$
\end{lemma}

\begin{proof}
The GCF with $a_n = -n!$ and $b_n = 1$ is formally associated
(via Euler's connection between continued fractions and power series)
to the alternating factorial series $f(x) = \sum_{n \ge 0} (-1)^n\,n!\,x^n$.
Its Borel transform cancels the factorial growth:
\[
  \mathcal{B}[f](t)
  \;=\;\sum_{n \ge 0}\frac{(-1)^n\,n!}{n!}\,t^n
  \;=\;\sum_{n \ge 0}(-t)^n
  \;=\;\frac{1}{1+t},
\]
in the canonical $k = 1$ normalization.

Multiplying all partial denominators by~$k$ is an equivalence
transformation of the continued fraction (Perron~\cite{Perron1957}).
In the Borel plane this rescales the Stieltjes kernel:
the substitution $t \mapsto t/k$ moves the simple pole from
$t = -1$ to $t = -k$, giving
\[
  \mathcal{B}_k(t) \;=\; \frac{k}{k+t}.
\]

Laplace (Borel) resummation along $[0,\infty)$ then yields
\[
  V(k)
  \;=\;\int_{0}^{\infty} e^{-t}\,\mathcal{B}_k(t)\,dt
  \;=\;\int_{0}^{\infty}\frac{k\,e^{-t}}{k+t}\,dt.
\]
The standard identity
$\displaystyle\int_0^\infty \frac{e^{-t}}{x+t}\,dt = e^{x}\,E_1(x)$
(valid for $\operatorname{Re} x > 0$; see~\cite[Ch.\,5]{DLMF})
with $x = k$ completes the proof:
\[
  V(k) = k\,e^{k}\,E_1(k). \qedhere
\]
\end{proof}

\paragraph{Numeric verification.}
A backward-recurrence computation of the GCF convergents at depth
$N = 200$ with $120$-digit precision matches the closed form to all
displayed digits.  For $k = 1,2,3$:
\[
\begin{array}{c|l|l}
k & V(k) = k\,e^k\,E_1(k) & \text{50-digit verification} \\
\hline
1 & 0.59634736232319407434\ldots & \checkmark \\
2 & 0.72265948505688951953\ldots & \checkmark \\
3 & 0.78606692861534416385\ldots & \checkmark
\end{array}
\]


% ── Bibliography entries (add to your .bib or \begin{thebibliography}) ──
%
% @book{Perron1957,
%   author    = {Perron, Oskar},
%   title     = {Die Lehre von den Kettenbr\"uchen},
%   edition   = {3rd},
%   publisher = {Teubner},
%   address   = {Stuttgart},
%   year      = {1957},
%   note      = {Equivalence transformations: \S\S\,1--4}
% }
%
% @book{Hardy1949,
%   author    = {Hardy, G. H.},
%   title     = {Divergent Series},
%   publisher = {Clarendon Press},
%   address   = {Oxford},
%   year      = {1949},
%   note      = {Borel summability: Ch.\,VIII}
% }
%
% @misc{DLMF,
%   title     = {{NIST} Digital Library of Mathematical Functions},
%   howpublished = {\url{https://dlmf.nist.gov/}},
%   note      = {Exponential integral $E_1$: \S\,6.2}
% }
%
% @book{Stieltjes1894,
%   author    = {Stieltjes, Thomas Jan},
%   title     = {Recherches sur les fractions continues},
%   journal   = {Annales de la Facult\'e des Sciences de Toulouse},
%   volume    = {8},
%   pages     = {J1--J122},
%   year      = {1894},
%   note      = {Stieltjes transform and moment representation}
% }
%
% @book{LorentzenWaadeland2008,
%   author    = {Lorentzen, Lisa and Waadeland, Haakon},
%   title     = {Continued Fractions: Convergence Theory},
%   edition   = {2nd},
%   publisher = {Atlantis Press},
%   year      = {2008},
%   note      = {Stern--Stolz theorem: Theorem~3.21}
% }

%
% Required packages (add to preamble if not already present):
%   \usepackage{amsmath,amssymb,amsthm}
%   \newtheorem{lemma}{Lemma}

\begin{lemma}[$k$-Shift Borel Regularization Identity]\label{lem:borel-gcf}
Let the generalized continued fraction (GCF) be defined by partial numerators
$a_n = -n!$ and constant partial denominators $b_n = k$ ($k > 0$),
with~$b_0 = k$.  The formal continued fraction is classically divergent
(by the Stern--Stolz theorem, since $|a_n/b_n| \to \infty$),
but its Borel regularization exists and equals
\[
  V(k)
  \;=\;\int_{0}^{\infty}\frac{k\,e^{-t}}{k+t}\,dt
  \;=\; k\,e^{k}\,E_1(k),
\]
where\/ $E_1(z) = \displaystyle\int_z^{\infty}\frac{e^{-t}}{t}\,dt$
is the exponential integral.
For $k=2$ we obtain $V(2) = 2e^{2}E_1(2) \approx 0.72266\ldots$
\end{lemma}

\begin{proof}
The GCF with $a_n = -n!$ and $b_n = 1$ is formally associated
(via Euler's connection between continued fractions and power series)
to the alternating factorial series $f(x) = \sum_{n \ge 0} (-1)^n\,n!\,x^n$.
Its Borel transform cancels the factorial growth:
\[
  \mathcal{B}[f](t)
  \;=\;\sum_{n \ge 0}\frac{(-1)^n\,n!}{n!}\,t^n
  \;=\;\sum_{n \ge 0}(-t)^n
  \;=\;\frac{1}{1+t},
\]
in the canonical $k = 1$ normalization.

Multiplying all partial denominators by~$k$ is an equivalence
transformation of the continued fraction (Perron~\cite{Perron1957}).
In the Borel plane this rescales the Stieltjes kernel:
the substitution $t \mapsto t/k$ moves the simple pole from
$t = -1$ to $t = -k$, giving
\[
  \mathcal{B}_k(t) \;=\; \frac{k}{k+t}.
\]

Laplace (Borel) resummation along $[0,\infty)$ then yields
\[
  V(k)
  \;=\;\int_{0}^{\infty} e^{-t}\,\mathcal{B}_k(t)\,dt
  \;=\;\int_{0}^{\infty}\frac{k\,e^{-t}}{k+t}\,dt.
\]
The standard identity
$\displaystyle\int_0^\infty \frac{e^{-t}}{x+t}\,dt = e^{x}\,E_1(x)$
(valid for $\operatorname{Re} x > 0$; see~\cite[Ch.\,5]{DLMF})
with $x = k$ completes the proof:
\[
  V(k) = k\,e^{k}\,E_1(k). \qedhere
\]
\end{proof}

\paragraph{Numeric verification.}
A backward-recurrence computation of the GCF convergents at depth
$N = 200$ with $120$-digit precision matches the closed form to all
displayed digits.  For $k = 1,2,3$:
\[
\begin{array}{c|l|l}
k & V(k) = k\,e^k\,E_1(k) & \text{50-digit verification} \\
\hline
1 & 0.59634736232319407434\ldots & \checkmark \\
2 & 0.72265948505688951953\ldots & \checkmark \\
3 & 0.78606692861534416385\ldots & \checkmark
\end{array}
\]


% ── Bibliography entries (add to your .bib or \begin{thebibliography}) ──
%
% @book{Perron1957,
%   author    = {Perron, Oskar},
%   title     = {Die Lehre von den Kettenbr\"uchen},
%   edition   = {3rd},
%   publisher = {Teubner},
%   address   = {Stuttgart},
%   year      = {1957},
%   note      = {Equivalence transformations: \S\S\,1--4}
% }
%
% @book{Hardy1949,
%   author    = {Hardy, G. H.},
%   title     = {Divergent Series},
%   publisher = {Clarendon Press},
%   address   = {Oxford},
%   year      = {1949},
%   note      = {Borel summability: Ch.\,VIII}
% }
%
% @misc{DLMF,
%   title     = {{NIST} Digital Library of Mathematical Functions},
%   howpublished = {\url{https://dlmf.nist.gov/}},
%   note      = {Exponential integral $E_1$: \S\,6.2}
% }
%
% @book{Stieltjes1894,
%   author    = {Stieltjes, Thomas Jan},
%   title     = {Recherches sur les fractions continues},
%   journal   = {Annales de la Facult\'e des Sciences de Toulouse},
%   volume    = {8},
%   pages     = {J1--J122},
%   year      = {1894},
%   note      = {Stieltjes transform and moment representation}
% }
%
% @book{LorentzenWaadeland2008,
%   author    = {Lorentzen, Lisa and Waadeland, Haakon},
%   title     = {Continued Fractions: Convergence Theory},
%   edition   = {2nd},
%   publisher = {Atlantis Press},
%   year      = {2008},
%   note      = {Stern--Stolz theorem: Theorem~3.21}
% }
